\documentclass[11pt]{article}
\usepackage[margin=0.85in]{geometry}
\usepackage{amsmath,amssymb}
\title{A cubic family of off-diagonal Vinogradov solutions}
\author{The Last Math Competition, conjecture 00000004174}
\date{}
\begin{document}
\maketitle

For the standard Vinogradov system with degree \(k=2\) and three variables
on each side, the off-diagonal count need not lose even a positive power
of the size bound compared with the diagonal count. This disproves the
conjecture's order-drop clause.

Let \(x,y\in\{1,\ldots,X\}^{3}\) be ordered triples. The two equations are
\[
 x_1+x_2+x_3=y_1+y_2+y_3,
 \qquad x_1^2+x_2^2+x_3^2=y_1^2+y_2^2+y_3^2.
\]
We call a solution diagonal if \(y\) is a permutation of \(x\), with an
arbitrary permutation allowed even when some coordinates repeat. The
constructed pairs satisfy both equalities and hence the corresponding
congruences modulo every positive modulus. Conversely, for modulus
\(q>3X^2\), the two congruences on this box imply the equalities because
the differences of the two sums have absolute value less than \(q\).

\paragraph{Construction.}
For any integer \(M\ge 1\), choose \(a,b,c\) independently from the intervals
\[
 M\le a,b\le 2M-1,\qquad 3M\le c\le 4M-1.
\]
Set \(t=a+b-2c\), \(w=(1,1,-2)\), and define
\[
 x=(3a,3b,3c),\qquad
 y=x-tw=(2a-b+2c,-a+2b+2c,2a+2b-c).
\]
The entries of \(w\) sum to zero, so subtracting \(tw\) preserves the first
power sum. Moreover \(x\cdot w=3t\) and \(w\cdot w=6\), giving
\[
 \|y\|^2-\|x\|^2=-2t(3t)+6t^2=0.
\]
All six coordinates lie in \([1,12M]\). Indeed,
\(3M\le x_1,x_2\le 6M-3\), \(9M\le x_3\le 12M-3\),
\(6M+1\le y_1,y_2\le 11M-4\), and \(1\le y_3\le 5M-4\).
Only \(x_3\) exceeds \(6M\) among the entries of \(x\), while both \(y_1\)
and \(y_2\) exceed \(6M\). Thus the triples cannot be permutations of one
another. Finally, the entries of \(x\) recover \(a,b,c\), so the construction
produces \(M^3\) distinct ordered pairs.

\paragraph{Counting consequence.}
For \(X\ge 24\), take \(M=\lfloor X/12\rfloor\ge X/24\). The off-diagonal
count is at least \(\lfloor X/12\rfloor^3\), while the diagonal count is at
most \(6X^3\) because each left triple has at most six reorderings. The
ratio of the off-diagonal count to the diagonal count is therefore bounded
below by a positive constant. The sharper classical asymptotic for the total
count, \(J_{3,2}(X)\sim (18/\pi^2)X^3\log X\) \cite{bb}, implies that the
off-diagonal count even has the leading asymptotic of the total count.

\paragraph{Formal verification.}
In \texttt{Vinogradov.lean}, the theorem \texttt{cubic\_off\_diagonal}
constructs an injection from \(\operatorname{Fin}(M)^3\) into the bounded
off-diagonal solutions at \(X=12M\). The codomain requires the two
equalities and both congruences modulo every modulus. The construction,
including bounds and exclusion of every permutation, is checked in Lean 4
without external dependencies.

\begin{thebibliography}{1}
\bibitem{bb}
V.~Blomer and J.~Br\"udern, \emph{The number of integer points on
Vinogradov's quadric}, Monatsh. Math. \textbf{160} (2010), 243--256.
\end{thebibliography}
\end{document}
